% !PRINTED 136
\weberpage{136}
%<<CONT

of the greatest common divisor is applied to these two functions according to the principles developed in the first section, the condition for common factors is obtained by setting the last constant remainder equal to zero, in the form of an equation between the coefficients,
\[
 F(a_1,a_2,\ldots,a_n)=0,
\]
where $F$ is an integral rational function of the arguments $a_1,a_2,
\ldots,a_n$, which is called the \gesperrt{discriminant} of $f$, and whose properties and mode of formation will be treated more fully in the next section. In any case, since equations of degree $n$ without multiple roots do exist at all, $F$ cannot vanish identically.

Consider now
\[
 F(a_1+\varepsilon_1,a_2+\varepsilon_2,
\ldots,a_n+\varepsilon_n).
\]
This function too, if $a_1,a_2,
\ldots,a_n$ are definite and $\varepsilon_1,\varepsilon_2,
\ldots,\varepsilon_n$ are indeterminate numbers, will not vanish identically; and one can choose the $\varepsilon$'s so that its value comes out different from zero, even when any arbitrary upper bound has been fixed for the absolute values of the $\varepsilon$'s.\footnote[1]{This will only later be rigorously justified. For the time being we presuppose it in these preliminary remarks, on which no further conclusions will be founded.}

One therefore sees from these considerations how an $a$-fold root $\alpha$ of $f(x)$, under continuous variation of the coefficients, radiates apart into $a$ simple roots, so that from this point of view too one is justified in considering the $a$-fold root as having arisen from the coincidence of $a$ simple roots.

If we do not assume the coefficient of the highest power of $x$ equal to $1$, and therefore examine the equation in the form
\begin{equation}
 a_0x^n+a_1x^{n-1}+a_2x^{n-2}+\cdots+a_{n-1}x+a_n=0,
\tag{21}
\end{equation}
then, if $a_n=0$, one or more of the roots are equal to zero, and equation (21) can be divided by $x$ or by a higher power of $x$, whereby an equation of correspondingly lower degree arises in which the last term is different from zero. We shall therefore assume from the outset that $a_n$ is different from zero, and now put in (21)
\begin{equation}
 x=\frac1y.
\tag{22}
\end{equation}
%CONT>>
